\subsection{A conditional cuspidal criterion for point separation}

\begin{claim}{AG.40}{Proved: a conditional cuspidal separation criterion}
\label{clm:AGcuspidalcriterion}
Let \(k\) be a finite field of characteristic \(p\), let
\(C/k\) be smooth, proper and geometrically connected of
genus at least two, and let \(P,Q\in C(k)\).
Write \(\Pi_C\) for the arithmetic fundamental group with
maximal geometrically prime-to-\(p\) quotient. Fix a prime
\(r\ne p\). Let \(G_P\) and \(G_Q\) be the maximal
cuspidally pro-\(r\) quotients of the corresponding
arithmetic prime-to-\(p\) groups of \(C-P\) and \(C-Q\),
relative to their maps to \(\Pi_C\).

Suppose there is a continuous isomorphism
\[
                  \beta:G_P\xrightarrow{\sim}G_Q
\]
over \(\operatorname{id}_{\Pi_C}\) which carries the
cuspidal decomposition group at \(P\) to a conjugate of
the cuspidal decomposition group at \(Q\). Then \(P=Q\).
The isomorphism is in particular over the identity of
\(G_k\), so it preserves arithmetic Frobenius.
\end{claim}

This is an implication from the existence of the specified
isomorphism. Its proof does not use a configuration-space
fiber exactness theorem. The required construction from
conjugate proper point sections is supplied separately by
Claims~\ref{clm:AGfixedrconfiguration} and
\ref{clm:AGcuspidaltransport}.

\begin{proof}
Assume \(P\ne Q\). We give the argument in five steps.

\emph{1. The point difference is \(p\)-primary.}
Projection of the cuspidal decomposition groups to the
proper quotient shows that the sections of \(P\) and
\(Q\) in \(\Pi_C\) are conjugate. The conjugator
may be made geometric: multiply it by an element of
the target point decomposition group with the same
arithmetic projection, which does not change the
conjugated subgroup. Thus these sections are
geometrically conjugate. Passing
to the geometric abelian quotient and the Jacobian
\(J_C\), their difference has zero Kummer class for
every prime \(\ell\ne p\). The Kummer map on the
\(\ell\)-primary part of \(J_C(k)\) is injective:
the usual Kummer exact sequences, followed by inverse
limit, give the injection of the \(\ell\)-adic
completion of \(J_C(k)\) into
\(H^1(G_k,T_\ell J_C)\). Since \(J_C(k)\) is finite,
this completion is precisely its \(\ell\)-primary part.
Consequently \([P-Q]\) has \(p\)-power order. There are
an integer \(e\geq0\) and \(f\in k(C)^\times\) with
\[
                  \operatorname{div}(f)=p^e(P-Q).
                                                        \tag{AG.40.1}
\]
Indeed this divisor is principal over \(\bar k\), and
its rational generator descends by Hilbert~90, since
the divisor is \(k\)-rational.

\emph{2. Every visible finite cover has its zeta function
and its number of missing points preserved.}
For the geometric group of \(C-P\), let \(K_P\) be
the kernel over the full proper geometric prime-to-\(p\)
group and let \(N_P=\ker(K_P\to K_P^{(r)})\).
Equation (AG.39.15) applies to every open subgroup
containing \(N_P\). Thus, for every finite cover
visible in \(G_P\), killing its cusp inertia and then
taking its maximal pro-\(r\) quotient gives the full
proper pro-\(r\) group of its smooth compactification.
This construction is natural for the arithmetic action.
The same holds for \(G_Q\).

The isomorphism \(\beta\) identifies the relevant
inertia groups, because it identifies the cuspidal
decomposition groups and their geometric kernels.
It therefore identifies the compactified
\(H^1(-,\mathbb Z_r)\) representations, including
Frobenius. The integral Weil polynomials, hence the
point counts of the compactifications over every
finite constant extension, agree. The missing points
also correspond: their decomposition groups and
residue degrees are given by the cuspidal
double-cosets and their arithmetic projections in
the respective open subgroups. It follows that
corresponding \emph{open} covers have the same
number of rational points over their constant
fields and over every further finite extension.

\emph{3. Noncuspidal point sections are transported.}
Extend constants temporarily so that a chosen
noncuspidal point \(R\) is rational. Let
\(D_R^P\subset G_P\) be its decomposition subgroup.
It maps isomorphically to the absolute Galois group
of the new constant field. Every open subgroup
\(H'\subset G_Q\) containing \(\beta(D_R^P)\)
has full arithmetic projection. The source cover
for \(\beta^{-1}(H')\) has a rational noncuspidal
point over \(R\), so the target open cover has a
rational noncuspidal point by Step~2.

The finite sets of rational points of these target
covers form a directed inverse system with the
finite intersection property: any finite collection
of conditions can be checked in the cover for the
intersection of the corresponding open subgroups.
Compactness gives a compatible point of the
pro-cover corresponding to the closed subgroup
\(\beta(D_R^P)\). Choose a geometric lift in the
universal pro-cover defined by \(G_Q\). Its
stabilizer is the image in \(G_Q\) of an actual
target decomposition subgroup, and is contained in
\[
        \bigcap_{H'\supset\beta(D_R^P)}H'
                                      =\beta(D_R^P).
\]
Both groups map isomorphically onto the constant-field
Galois group, so this containment is equality.
Equivalently, \(\beta\) carries \(D_R^P\), up to
conjugacy, to the decomposition subgroup of an actual
noncuspidal target point.

For the full proper prime-to-\(p\) group, the
exceptional geometric point set is finite
\cite[Proposition~1.8(vi)]{AR_ST2009}.
Outside that set, a proper decomposition subgroup
has a unique point label. Since \(\beta\) is over
the identity of \(\Pi_C\), the target point in
Step~3 is therefore \(R\) itself whenever \(R\)
is outside this finite set and \(\{P,Q\}\).
This statement applies to every closed point after
its finite constant extension.

\emph{4. A difference cocycle turns point transport
into constant Kummer evaluations.}
Let \(M=\mathbb Z_r(1)\). Denote by
\[
  1\longrightarrow M\longrightarrow E_P
     \longrightarrow\Pi_C\longrightarrow1,
  \qquad
  1\longrightarrow M\longrightarrow E_Q
     \longrightarrow\Pi_C\longrightarrow1
                                                        \tag{AG.40.2}
\]
the maximal geometrically cuspidally central
quotients of \(G_P,G_Q\). Their geometric
extension classes are the degree-one cycle classes
of \(P\) and \(Q\) in
\(H^2(\Delta_C,M)\simeq\mathbb Z_r\).
Here is a construction avoiding any assertion
about a completed configuration or torsor fiber.
Use the Kummer cycle class of
\(\mathcal O_C(P)\) in
\(H^2(\Pi_C,M)\) to define the extension.
Its restriction to \(C-P\) is canonically
split by the canonical section of that line
bundle, giving a homomorphism from the
punctured-curve group to the extension.
The local residue maps the puncture inertia
onto the canonical generator of \(M\).
Conversely, the maximal geometrically central
cuspidal kernel is generated topologically by
this one inertia image: the full cuspidal
kernel is normally generated by the puncture
inertia, and geometric conjugation has been
killed. Hence the displayed homomorphism
identifies this central kernel with \(M\).
The same construction applies to \(Q\).
The geometric degree can also be checked directly
from the punctured surface presentation
\cite[XIII, Corollary~2.12]{AR_SGA1}: the puncture
word is the product of the \(g\) handle commutators.
After geometric cuspidal centralization it is one central
generator. Its image has infinite \(r\)-adic order:
send the first handle generators to \(1+e_{12}\)
and \(1+e_{23}\) in the upper-unitriangular
Heisenberg group over \(\mathbb Z_r\), and the
remaining handle generators to one. The puncture word
maps to \(1+e_{13}\). The resulting central relator
is the standard degree-one orientation extension.
This verifies the geometric degree-one class directly,
in addition to the cycle-class construction.
These are one-dimensional Kummer and localization
constructions, with no configuration-space fiber
injectivity input.

The isomorphism \(\beta\) induces an extension
isomorphism \(\bar\beta:E_P\to E_Q\).
Its multiplier on the kernel \(M\) is one:
it is over the identity on \(\Delta_C\), and
both degree-one extension classes are the same
generator of \(H^2(\Delta_C,M)\).
On the fiber product \(E_P\times_{\Pi_C}E_Q\)
define
\[
                 \chi(u,v)=\bar\beta(u)v^{-1}\in M.
                                                        \tag{AG.40.3}
\]
This is a continuous crossed homomorphism for the
cyclotomic action. Pull it back along the natural
map from the arithmetic prime-to-\(p\) group of
\(C-\{P,Q\}\) to that fiber product. Its residues
at \(P,Q\) are respectively \(+1,-1\), with
the canonical Kummer orientations. Therefore
\[
                         \kappa_r(f)-p^e\chi
                                                        \tag{AG.40.4}
\]
has zero residues, kills both inertia groups, and
descends to a class in \(H^1(\Pi_C,M)\).

This last cohomology group consists of constant
Kummer classes. For completeness, the
Hochschild--Serre exact sequence gives
\[
 H^1(G_k,M)\xrightarrow{\sim}H^1(\Pi_C,M),
                                                        \tag{AG.40.5}
\]
because \(H^1(\Delta_C,M)^{G_k}=0\): its
Frobenius eigenvalues have nonzero Weil weight,
so none is one. Finite-field
Kummer theory represents the class of (AG.40.4) by
some \(a\in k^\times\).

We check that the evaluation of \(\chi\) at
each point \(R\) from Step~3 is zero as a
cohomology class. Work over its finite residue
field so that \(R\) is rational. The target
proper point section satisfies
\[
                       N_{\Pi_C}(D_R)=D_R.
                                                        \tag{AG.40.6}
\]
Indeed a normalizer element can be divided by
the unique element of \(D_R\) with the same
arithmetic projection; the quotient is geometric
and centralizes \(D_R\), hence is trivial by
the standard Frobenius-weight centralizer
argument (equivalently
\cite[Lemma~1.7]{AR_ST2009}).

Choose a conjugator between
\(\beta(D_R^P)\) and the target subgroup
\(D_R^Q\). Its projection to \(\Pi_C\)
lies in the normalizer in (AG.40.6). Remove that
projection by multiplying by the corresponding
element of \(D_R^Q\). The resulting conjugator
lies in the cuspidal kernel. Its image in
\(E_Q\) lies in \(M\); conjugation by it
changes a section only by a coboundary.
Consequently (AG.40.3) restricts to a coboundary
on the actual point section of \(R\).
Evaluating (AG.40.4) thus gives
\[
       \kappa_{r,k(R)}\bigl(f(R)/a\bigr)=0
                                                        \tag{AG.40.7}
\]
for every closed \(R\) outside the fixed finite
exceptional set and \(\{P,Q\}\).

\emph{5. Constant Kummer evaluations force a
constant function.}
If \(f\) were nonconstant, its map from
\(C_{\bar k}\) to \(\mathbb P^1_{\bar k}\)
would be surjective. Choose a nontrivial
\(r\)-power root of unity \(\zeta\) such that
\(a\zeta\) avoids the finitely many images
of the excluded points. There are infinitely
many choices because \(r\ne p\).
Take a closed point \(R\) with
\(f(R)=a\zeta\). Its finite residue field
contains \(\zeta\), whose nontrivial
\(r\)-primary Kummer class in that field
contradicts (AG.40.7). Thus \(f\) is constant,
contradicting (AG.40.1) for \(P\ne Q\).
This proves the conditional criterion.
\end{proof}

\paragraph{The separate existence construction and source scope.}
Claims~\ref{clm:AGfixedrconfiguration} and
\ref{clm:AGcuspidaltransport} below supply the required
cuspidal isomorphism from conjugate proper point sections,
using a direct one-prime relative configuration argument.
They do not use the full prime-to-\(p\) configuration
fiber assertion excluded by AG39.
For clarity about the earlier literature route,
Saidi--Tamagawa's Corollary~2.3 constructs the compatible
cuspidally abelian comparison only when the chosen cusp
coordinates lie outside its exceptional sets; its
Theorem~2.7 refers to that comparison. The full
configuration-fiber assertion in the cited Mochizuki
author versions requires a quotient-convention or
source-scope reconciliation with AG39. Neither an
unverified exceptional-coordinate extension of those
statements nor that full configuration assertion is an
input to AG40--AG42.
